% !PRINTED 512
\weberpage{512}

%<<CONT
\gesperrt{geraden Teiler und so, daß, wenn $x$ gerade ist, $y\equiv 2\pmod 4$, ferner so, daß
\[
n=\tfrac14(y^2-Dx^2)
\]
eine ungerade Quadratzahl ist. Es sei ferner $M$ durch die Formel (12) bestimmt. Dann hat die Funktion $H_{-D}(u)$ einen Teiler
\weq{22}{\Phi(M,u),}
der mit keiner der Funktionen
\weq{23}{\Phi(-M,u),\quad \Phi(iM,u),\quad \Phi(-i,Mu)}
% [sic] printed third term $\Phi(-i,Mu)$ misplaces the comma; the sign orbit requires $\Phi(-iM,u)$
eine gemeinsame Wurzel hat.}

\begin{center}
§ 137.\quad \gesperrt{Zurückführung ungerader Diskriminanten auf gerade.}
\end{center}

Nach einem von \gesperrt{Kronecker} ausgesprochenen Satz läßt sich die Klassengleichung unter Adjunktion von Quadratwurzeln in so viele Faktoren zerlegen, als es für die betreffende Klassengleichung Geschlechter der Formenklassen gibt, und zwar so, daß jedem dieser Faktoren nur die Klasseninvarianten genügen, für die die entsprechenden quadratischen Formen $(A,B,C)$ zu einem und demselben Geschlecht gehören. Diese Zerlegung erhalten wir aus den im vorigen Paragraphen bewiesenen Satz 1. % [sic] printed "aus den ... Satz"; grammatically "aus dem"

Daß dieser Satz in der Form, in der wir ihn ausgesprochen haben, sich nur auf gerade Diskriminanten $D$ bezieht, ist für den Beweis des \gesperrt{Kronecker}schen Satzes keine Beschränkung.

Denn wenn $D\equiv 1\pmod 8$ ist, dann ist nach § 123 der Klassenkörper $\mathfrak{K}(D)$ identisch mit dem Klassenkörper $\mathfrak{K}(4D)$. Ist aber $D\equiv 5\pmod 8$, dann ist der Grad von $\mathfrak{K}(4D)$ dreimal so groß als der von $\mathfrak{K}(D)$, aber der letztere ist in dem ersteren enthalten, und die Anzahl der Geschlechter für $\mathfrak{K}(D)$ und $\mathfrak{K}(4D)$ ist die gleiche (§ 104).

Ist $u=j(\omega)$ eine Klasseninvariante für $D\equiv 1\pmod 8$, so sind
\weq{1}{v=j(\omega')=j(2\omega),\quad j\left(\dfrac{\omega}{2}\right),\quad j\left(\dfrac{\omega+1}{2}\right)}
Klasseninvarianten von $4D$, und $u$ ist eine rationale Funktion von $v$:
\weq{2}{u=\varphi(v),}
die ungeändert bleibt, welche der drei Größen (2) man auch für $v$ setzen mag. % [sic] printed reference (2) should be (1), which contains the three quantities
Alle Charaktere der Formen, deren Wurzeln $\omega$ und $\omega'$ sind, haben denselben Wert.
